% Generated draft: requires mathematical and source review.
\documentclass[11pt]{article}
\usepackage[margin=25mm]{geometry}
\usepackage{fontspec}
\setmainfont{Latin Modern Roman}
\usepackage{amsmath,amssymb,mathtools,unicode-math}
\setmathfont{Latin Modern Math}
\usepackage{enumitem,needspace,xurl,hyperref}
\setlength{\parindent}{0pt}
\setlength{\parskip}{5pt}
\setlength{\emergencystretch}{3em}
\newcommand{\N}{\mathbb{N}}
\newcommand{\Z}{\mathbb{Z}}
\newcommand{\Q}{\mathbb{Q}}
\newcommand{\R}{\mathbb{R}}
\newcommand{\C}{\mathbb{C}}
\newcommand{\F}{\mathbb{F}}
\newcommand{\A}{\mathbb{A}}
\newcommand{\PP}{\mathbb{P}}
\newcommand{\E}{\mathbb{E}}
\newcommand{\eps}{\varepsilon}
\newcommand{\dd}{\,\mathrm{d}}
\begin{document}
\section*{UnsolvedMath ID: 30006319}
\textbf{Source problem code:} catalogue entry
\subsection{Definitions and mathematical statement}
The supplied record specifies a conjectural template whose full hypotheses and categorical conventions are not supplied. Let \((G,X)\) denote the hyperspherical pair mentioned in the record, with \(G\) acting on \(X\), and let \(G^\vee\) and \(M^\vee\) denote its dual group and dual Hamiltonian space. The definition of a hyperspherical pair, the permitted base fields, and the construction of this dual data are not included in the record and remain prerequisites for a complete statement.

Write \(\mathcal D\) for the formal disk and \(\mathring{\mathcal D}\) for its punctured disk, in the unspecified geometric setting of the source. The notation \(G(\mathcal D)\) denotes the corresponding group of disk-valued points, acting on \(X(\mathring{\mathcal D})\). Write \(\operatorname{Shv}_{G(\mathcal D)}(X(\mathring{\mathcal D}))\) for the source's category of equivariant sheaves, whose sheaf theory, coefficients, and derived conventions are unspecified. Write \(\operatorname{QCoh}(M^\vee)^{G^\vee}\) for the category of \(G^\vee\)-equivariant quasi-coherent sheaves on the dual Hamiltonian space, subject to the unresolved categorical decorations in the original extraction. Let \(\mathcal H_G\) denote the spherical Hecke category referred to by the source.

\textbf{Conjectural statement, with unresolved qualifications.} For every hyperspherical pair satisfying the additional conditions referenced by the source as \([3,\text{ Conjecture }7.5.1]\), there should exist an equivalence
\[
\Phi:\operatorname{Shv}_{G(\mathcal D)}\!\bigl(X(\mathring{\mathcal D})\bigr)
\xrightarrow{\ \sim\ }
\operatorname{QCoh}(M^\vee)^{G^\vee}
\]
of \(\mathcal H_G\)-module categories. Thus, for every \(A\in\mathcal H_G\) and every object \(\mathcal F\) of the left-hand category, the equivalence should intertwine the specified module actions through coherent isomorphisms
\[
\Phi(A\star\mathcal F)\simeq A\cdot\Phi(\mathcal F).
\]
The actions themselves and the coherence conventions must be taken from the intended source framework; they are not defined by the supplied record. The original extraction also indicates a spectral description of \(\mathcal H_G\) involving \(\operatorname{QCoh}(\check{\mathfrak g}^{*})^{G^\vee}\), but contains unreadable decorations. No precise graded, shifted, or otherwise modified version of that description can be recovered from this record alone.

The source's accompanying remark introduces a vector space \(V_X\) and a subgroup \(G_X^\vee\subseteq G^\vee\), and describes
\[
M^\vee=G^\vee\times_{G_X^\vee}V_X.
\]
Here the associated-space notation requires the intended \(G_X^\vee\)-action on \(V_X\). In a setting where this presentation and the relevant quotient stacks are defined, the remark gives
\[
\begin{aligned}
\operatorname{QCoh}(M^\vee)^{G^\vee}
&\simeq\operatorname{QCoh}([M^\vee/G^\vee])\\
&\simeq\operatorname{QCoh}([V_X/G_X^\vee])\\
&\simeq\operatorname{QCoh}(V_X)^{G_X^\vee}.
\end{aligned}
\]
Consequently the source also proposes this last category as an equivalent target, carrying the transported Hecke-module structure and all the same unresolved qualifications. The missing conditions in \([3,\text{ Conjecture }7.5.1]\), categorical conventions, and definitions prevent this template from being a fully specified standalone conjecture.
\subsection{Short English statement}
The proposed local relative Langlands equivalence identifies equivariant sheaves on punctured-disk points of a hyperspherical space with equivariant quasi-coherent sheaves on its dual Hamiltonian space, compatibly with the spherical Hecke action. The source also expresses the target using a vector space \(V_X\) and a subgroup \(G_X^\vee\). Essential hypotheses, sheaf conventions, and conditions referenced elsewhere remain unspecified, so this is a qualified statement-only draft.
\subsection{Sources}
\begin{itemize}
\item Supplied local UnsolvedMath registry record, permanent ID \(30006319\), source problem code \texttt{catalogue entry}, title: \emph{Dual Hamiltonian Geometry of Equivariant Sheaf Categories}. The record belongs to the Oberwolfach Reports Open Problems collection and gives proposed year 2025; no individual proposer is identified.
\item \emph{Arbeitsgemeinschaft: Relative Langlands Duality} (2025), Oberwolfach Reports, DOI 10.4171/owr/2025/18, pp. 811--882. Supplied report citation and link: \url{https://doi.org/10.4171/owr/2025/18}.
\item Ben-Zvi, Sakellaridis, and Venkatesh, \emph{Relative Langlands Duality}. Supplied literature link: \url{https://www.math.ias.edu/~akshay/research/BZSVpaperV1.pdf}.
\item Devalapurkar, \emph{Derived geometric Satake for PGL products}. Supplied literature link: \url{https://sanathdevalpurkar.github.io/files/PGL2-cubes.pdf}.
\item The original extraction refers to additional conditions in \([3,\text{ Conjecture }7.5.1]\). The supplied record does not provide the numbered bibliography needed to identify reference \([3]\) reliably.
\item Literature assessment date: \texttt{2026-08-22}. This is supplied dated curation, not independently verified here. The supplied assessment labels the problem open, describes it as the general local conjecture, and says that only special cases are known. The background explicitly requests independent verification of time-sensitive status and cited claims.
\item The registry labels the statement \texttt{corrected\_verified} and reports that the official report was checked and its notation repaired. These are supplied source-verification claims, not verification performed for this draft.
\end{itemize}
\end{document}
